\begin{abstract}
\label{sec:abstract}
% What problem; why important
% Existing multi-LLM frameworks leverage the complementary expertise of different models, achieving stronger capabilities than scaling a single model.
Multi-LLM systems harness the complementary strengths of diverse Large Language Models, achieving performance and efficiency gains that are not attainable by a single model.
% Current solution
In existing designs, LLMs communicate through text, forcing internal representations to be transformed into output token sequences.
This process both loses rich semantic information and incurs token-by-token generation latency. 
Motivated by these limitations, we ask: \emph{Can LLMs communicate beyond text?} 
% Oracle?
Oracle experiments show that enriching the KV-Cache semantics can improve response quality without increasing cache size, supporting KV-Cache as an effective medium for inter-model communication.
% we identify it as an ideal communication medium with complementary perspectives of different LLMs on the same input.
% Ours
Thus, we propose Cache-to-Cache (\name), a new paradigm for direct semantic communication between LLMs.
% Method
\name uses a neural network to project and fuse the source model’s KV-cache with that of the target model to enable direct semantic transfer. 
A learnable gating mechanism selects the target layers that benefit from cache communication.
Compared with text communication, \name utilizes the deep, specialized semantics from both models, while avoiding explicit intermediate text generation.
% Experiments
Experiments show that \name achieves 6.4-14.2\% higher average accuracy than individual models. 
It further outperforms the text communication paradigm by approximately 3.1-5.4\%, while delivering an average 2.5× speedup in latency.
Our code is available at \url{https://github.com/thu-nics/C2C}.
\end{abstract}